% ============================================================================
%  C — file phần: Viết rõ
%  Ngày tạo: 2026-09-07 | Sửa gần nhất: 2026-09-07 | Ôn gần nhất: chưa
% ============================================================================
\documentclass[../english.tex]{subfiles}
\begin{document}
\part{VIẾT RÕ}
\begin{tomtat}
Phần C đảo chiều phần B: bây giờ bạn là người phải làm cho người khác hiểu đúng. Năm chương đi từ nhỏ tới lớn: câu (11), nối câu thành mạch (12), đoạn và bài (13), chọn đúng mức cam kết (14), và sửa bài (15). Nguyên tắc xuyên suốt: văn rõ không phải văn đơn giản mà là văn \emph{đặt thông tin đúng chỗ người đọc chờ nó}. Nếu chỉ nhớ một điều: người đọc quyết định nghĩa của câu dựa trên \emph{vị trí} của từ ngữ nhiều hơn bạn tưởng --- đầu câu là chỗ nối với cái đã biết, cuối câu là chỗ nhấn cái mới.
\end{tomtat}

Có một quan niệm sai rất phổ biến ở người viết tiếng Anh không phải bản ngữ: viết tốt là viết ``sang'' --- từ dài, câu phức, giọng bị động, danh từ hoá. Quan niệm này không phải tự nhiên mà có; nó đến từ việc quan sát các bài báo đã xuất bản và bắt chước bề mặt của chúng. Nhưng những người nghiên cứu văn phong khoa học đều đi tới cùng một kết luận ngược lại: các đặc điểm đó là \emph{tật} chứ không phải chuẩn, và những bài dễ đọc nhất trong mỗi ngành thường là những bài viết đơn giản nhất ở cấp câu.

Phần này dựa trên hai truyền thống. Truyền thống thứ nhất, gắn với Joseph Williams, xét câu ở cấp \emph{ai làm gì}: câu rõ khi chủ thể của hành động là chủ ngữ và hành động là động từ. Truyền thống thứ hai, gắn với George Gopen và Judith Swan, xét câu ở cấp \emph{người đọc chờ gì ở đâu}: mỗi vị trí trong câu có một chức năng, và văn khó đọc thường là văn đặt thông tin sai vị trí. Chương 11 và 12 lần lượt là hai truyền thống đó.

Ba chương còn lại mở rộng ra ngoài câu. Chương 13 nói về đơn vị lớn hơn --- đoạn, phần, và bài --- cùng những khuôn mà cộng đồng khoa học đã quen chờ đợi. Chương 14 xử lý một việc mà người viết không bản ngữ hay làm sai theo cả hai chiều: nói mạnh hơn bằng chứng cho phép, hoặc rào đón tới mức không còn khẳng định gì. Và chương 15 nói về phần công việc mà mọi người viết giỏi đều thừa nhận là quan trọng nhất nhưng ít ai dạy: sửa.
\subfile{00-map}
\subfile{11-cau-ro/cau-ro}
\subfile{12-mach-van/mach-van}
\subfile{13-doan-bai-tom-tat/doan-bai-tom-tat}
\subfile{14-noi-dung-muc/noi-dung-muc}
\subfile{15-sua-bai/sua-bai}
\ifSubfilesClassLoaded{\printbibliography[title={Nguồn}]}{}
\end{document}
